\begin{conclusions}
% R08 en ../write_thesis_roadmap.md. Una conclusión por objetivo específico, en el orden de
% Introduction.tex; ninguna afirma algo que no esté en los capítulos 1 a 5. Sin referencias nuevas.

Esta tesis se propuso desarrollar un generador híbrido de historias con LLMs, que combinara un plan
validado con la simulación de personajes con memoria propia, y una metodología para evaluarlo basada
en recuentos verificables y comparaciones humanas por pares a ciegas. A continuación se recoge lo
alcanzado en cada objetivo específico, en el orden en que se plantearon.

El análisis del estado del arte (Capítulo~\ref{chapter:state-of-the-art}) mostró que las tres
familias de generadores resuelven partes distintas del problema. Los dirigidos por la trama que se
revisaron garantizan la estructura y los compromisos, pero escriben a todos los personajes desde la
visión del autor; los que simulan personajes dan agencia, pero dejan la trama a la emergencia, y los
híbridos clásicos dependían de mundos formalizados a mano. En la literatura revisada, que no es
exhaustiva, no se encontró un generador con LLMs que reuniera un plan validado por código,
personajes cuya percepción calcula el código y un narrador que cura su registro sin inventar; los dos
más cercanos, Yu et al. y PLOTTER, reúnen parte de esas piezas. En la evaluación, la literatura revisada desaconseja las escalas
absolutas y los jueces LLM como medida de calidad, y respalda las comparaciones por pares.

El diseño del generador (Capítulo~\ref{chapter:generator}) separa qué pasa, cómo pasa y cómo se
cuenta. Un plan, validado por código como grafo acíclico y con un contrato de promesas anclado a sus
eventos, decide qué pasa; una función de actores que no ven el plan y recuerdan solo lo que
percibieron decide cómo pasa, con un director que sugiere y cierra cada \textit{beat}, y un narrador
que cura el registro sin inventar decide cómo se cuenta. La configuración de referencia mantiene
mecanismos opcionales que pueden activarse para estudiar su efecto sobre las medidas de la prosa y
la preferencia humana.

La metodología de evaluación (Capítulo~\ref{chapter:evaluation}) articula tres evaluaciones en las
que ninguna cifra automática se presenta como juicio de calidad: 24 rasgos núcleo, 18 de ellos
recuentos de un extractor con una cita que el código verifica literalmente y 6 medidas calculadas por
código, tres de ellas sobre lo que extrae el modelo; una evaluación humana por pares a ciegas con
tres preguntas y una línea base de un solo \textit{prompt}, y un juez que intenta reproducir el orden
humano a partir de los rasgos y se valida fuera de muestra, frente al azar y frente a la longitud
sola.

La implementación (Capítulo~\ref{chapter:implementation}) se organizó en cinco paquetes con
dependencias en un solo sentido, de modo que la evaluación no importa el generador y mide del mismo
modo cualquier historia guardada. Cada decisión del código vive en un validador puro con pruebas: el 1
de noviembre de 2026 la suite pasó 907 pruebas, entre ellas las que comprueban que ningún actor ve el
plan y que, con el inventario apagado, el generador queda como estaba. Cada ejecución queda en disco con su
versión y su huella, y las tablas del estudio, del juez y del ranking se regeneran desde esos
artefactos. La cuota diaria
del servicio de modelos fue una restricción práctica que el sistema administra, pero no evita.

La evaluación experimental (Capítulo~\ref{chapter:experiments}) documentó una configuración de
referencia y un conjunto horizontal de 24 ejecuciones: cinco bases simuladas, nueve líneas base,
dos narrativas desde el plan y ocho variantes de narración. La validación registrada del extractor
comprende doce historias, con repetición en seis; sus resultados se presentan por rasgo, con el
número de citas juzgadas y de pares efectivamente comparables. La generación terminada y la
extracción completa se verifican por separado.

El estudio humano, según la actualización del autor, reunió ocho graduados de Ciencia de la
Computación, incluido él mismo, con diferentes frecuencias de lectura. Cada participante comparó
cinco pares de diez textos distintos mediante Telegram y respondió tres preguntas por par, para un
total de 120 respuestas declaradas. Este alcance describe la participación; las preferencias por
criterio y el acuerdo requieren el registro individual, que no se sustituye por los votos del
piloto anterior. Los cinco pares disjuntos tampoco permiten identificar un ranking humano global.

Los contrastes responden a cuestiones diferentes. La matriz conserva escritura directa y simulación
vinculadas por plan; la preferencia entre ambas necesita relacionarse con los pares humanos. El autor
informa de resultados favorables a la memoria propia, cuya magnitud debe documentarse con las
ejecuciones y recuentos correspondientes. Las sesiones conservadas de mecanismos opcionales aportan
una observación compatible con la tercera hipótesis: reducir defectos medidos no aseguró que se
prefiriera el texto modificado. Extender estas observaciones a otras historias y lectores exige una
muestra mayor y la trazabilidad completa de las comparaciones.

Las \textbf{aportaciones} son el generador y la metodología de evaluación. El primero combina
planificación, simulación y narración con comprobaciones explícitas en sus interfaces. La segunda
articula recuentos citados, preferencias humanas por pares y un juez aprendido de esas preferencias.
El rango de acierto comunicado para el juez es de 0,50 a 0,65; sin su desglose y partición no se
interpreta como rendimiento verificado fuera de muestra. El corpus de 122 historias destinado a
ordenar se distingue de los textos que aportan votos al entrenamiento y a la validación.

Las \textbf{limitaciones} se refieren a la escala, la representatividad y la documentación de los
resultados. Los participantes comparten formación y el autor conoce el sistema; las lecturas
adicionales no cuentan como nuevos participantes. La cuota limita las ejecuciones y extracciones,
y los informes históricos no reemplazan los registros del estudio actualizado. Ampliar la muestra,
conectar los pares y conservar las respuestas y predicciones permitirá estimar la incertidumbre.
Usar el juez como función objetivo y buscar configuraciones con un algoritmo genético quedan como
trabajo futuro, una vez demostrada su generalización.

\end{conclusions}
